% SPDX-FileCopyrightText: 2026 Will Cook
% SPDX-License-Identifier: CC-BY-4.0
%
% One of the Erdős Problem Notes: a short standalone treatment of a single
% problem in the ErdosProblems expansion library.  Build with
% `tectonic erdos-1049-rational-base-lambert.tex` or `make`.
\documentclass[11pt]{article}
\input{problem-note-preamble}
\renewcommand{\commit}{99f4bf47422abbd8757cbb22b50ba079d764d3a7}
\renewcommand{\commitshort}{99f4bf47422a}
\hypersetup{
  pdftitle={Irrationality of F(31/4) and the exact normalized Hankel order (Erdős Problem 1049)},
  pdfsubject={Erdős Problem Notes: Problem 1049},
  pdfkeywords={irrationality, Lambert series, rational base, Hankel determinant, q-logarithm, Padé approximation, Lean 4},
}

\runninghead{Erd\H{o}s \#1049: an irrational base and the exact Hankel order}
\title{Irrationality of $F(31/4)$ and the\\Exact Normalized Hankel Order}
\subtitle{Cyclotomic cancellation and a unique least-order moment configuration}
\author{Will Cook}
\date{September 2026}

\begin{document}
\maketitle
\noteprovenance

\begin{abstract}
For coprime integers $a>b\ge1$, the Lambert value
$F(a/b)=\sum_{n\ge1}((a/b)^n-1)^{-1}$ is irrational whenever
$\log b/\log a<\theta^*$, where the explicit constant below satisfies
$\theta^*\approx0.4056830213840605$.
In particular, $F((31/4)^r)$ is irrational for every integer $r\ge1$.
The same forms bound its irrationality exponent uniformly in $r$.
This is an ordinary proof: it uses the polynomial conclusion of Zudilin's
Lemma~7 before the source's integer-specialisation step; Lean checks finite
subclaims, not the irrationality theorem.
Cyclotomic factors are cancelled before the rational-base denominator is
cleared; the degree of the resulting polynomial pair determines the cost.
The argument does not cover $3/2$.
For the normalised Hankel determinant in the second construction we prove
\[
 \operatorname{ord}_qV_N^*=\frac{N(N-1)(2N-1)}6,
 \qquad [q^{\operatorname{ord}_qV_N^*}]V_N^*
       =\frac{(N!)^2(N+1)!}{2^N}.
\]
The exponent and coefficient come from a unique least-order tuple in a formal
moment expansion.  A quantitative selector criterion separates local
divisibility from the additional estimates needed for a small nonzero real
remainder.
\end{abstract}

\section{Introduction}\label{sec:problem}

For $t>1$, expansion of each geometric series gives
\[
 F(t)=\sum_{n\ge1}\frac1{t^n-1}
     =\sum_{n\ge1}\frac{\tau(n)}{t^n},
\]
where $\tau(n)$ is the number of positive divisors of $n$.
Erd\H{o}s Problem~\#1049 asks for irrationality at every rational $t>1$
\cite[p.~102]{erdos1988}.  The theorem below proves an explicit sufficient
region and settles the power family in the title.

Cyclotomic cancellation produces positive forms $\Lambda_n=U_nF-V_n$
with integral coefficient polynomials of degree at most $W_n$, and
\[
 \log\!\bigl(b^{W_n}\Lambda_n(a/b)\bigr)
 =\bigl(C_1\log b-C_0\log a\bigr)n^2+o(n^2).
\]
A negative exponent yields positive integral linear forms tending to zero.

The Hankel argument concerns a different family.  Its moment expansion makes
one increasing index tuple responsible for the first coefficient.  The final
sections explain what local cancellation can establish at $3/2$ and specify
the real estimate that would complete that approach.  The complete catalogue
of other constructions and their failure witnesses is retained in the long
record.


\section{A rational base at which $F$ is irrational}
\label{sec:rational-base-irrationality}

Write $\psi_1(u)=\sum_{k\ge0}(k+u)^{-2}$.  Let $\mathcal I$ be the thirteen
intervals listed in the proof and put
\[
 C_1=\frac{1091}{2},\qquad
 J=\sum_{[u,v)\in\mathcal I}\bigl(\psi_1(u)-\psi_1(v)\bigr),\qquad
 C_0=266-\frac3{\pi^2}(225-J).
\]
Thus $\theta^*=C_0/C_1$ and $\mu=C_1/C_0$ are defined exactly.


\begin{theorem}[rational-base region]
\label{res:rational-base-threshold}
Let $a>b\ge1$ be coprime integers with
\[
\begin{gathered}
 b^{\mu}<a,\qquad\text{equivalently}\qquad
 \frac{\log b}{\log a}<\theta^*,\\
 \theta^*=\frac{C_0}{C_1}=0.4056830213840605\ldots,\\
 \mu=\frac{C_1}{C_0}=2.4649786835749750\ldots.
\end{gathered}
\]
Then $F(a/b)$ is irrational.
\end{theorem}

\noindent The cancellation is performed before homogenisation, so the
clearing degree is the degree of the cancelled pair.  The value $3/2$ is
outside this sufficient region: $J\le\psi_1(1/14)-\psi_1(1)<196$ gives
$\theta^*<266/(1091/2)<1/2<\log2/\log3$.

\begin{proof}[Proof of Theorem~\ref{res:rational-base-threshold}]
Set
\[
 a_0=14n+1,\quad a_1=12n+1,\quad a_2=14n+1,\quad\beta=27n+2,
 \qquad N=15n.
\]
Use the coefficient pair $A_n,B_n$ of the source identities
\cite[(8)--(11)]{zudilin2004}, with $H_n=A_nF-B_n$.
Put $D_N(X)=\prod_{\ell=1}^{N}\Phi_\ell(X)$,
$M_n=266n^2+34n+1$, and
\[
 \Omega_n(X)=\prod_{\ell=2}^{N}\Phi_\ell(X)^{\nu_\ell},
 \qquad \nu_\ell=\omega(n/\ell),
\]
where the periodic function
\[
\begin{split}
\omega(x)=\max\{0,&\ \lfloor14x\rfloor+\lfloor13x\rfloor
                  -\lfloor12x\rfloor-\lfloor15x\rfloor,\\
                &\ 2\lfloor14x\rfloor-\lfloor13x\rfloor
                  -\lfloor15x\rfloor\}
\end{split}
\]
is zero or one.  Its support in $[0,1)$ consists of
\[
\begin{gathered}
\,[1/14,1/12),\ [1/7,1/6),\ [3/14,1/4),\ [2/7,1/3),\\
[5/14,2/5),\ [3/7,7/15),\ [1/2,8/15),\ [4/7,3/5),\\
[9/14,2/3),\ [5/7,11/15),\ [11/14,4/5),\\
[6/7,13/15),\ [13/14,14/15).
\end{gathered}
\]
These are the intervals $\mathcal I$ used to define $J$.
The zero-one values and this thirteen-interval support are
\href{https://github.com/wcook04/plectis-erdos/blob/f36a98bf3d3e6f65f1074e3b800e3293d5b8a51a/ErdosProblems/Erdos1049/PaperOmegaIndicatorR7.lean#L1297}{Lean source}.

\paragraph{Integral polynomials and their degrees.}
The source's polynomial inclusion is Lemma~7, display~(23).
Its parameter vector is $n(13,14,12,14,15,13)$, its maximum is $15n$, and
$\beta-a_1-a_2=n>0$.  Its hypotheses therefore hold for every $n\ge1$.
It gives
\[
 \Lambda_n(X)=X^{-M_n}\frac{D_N(X)}{\Omega_n(X)}H_n(X)
            =U_n(X)F(X)-V_n(X),\qquad U_n,V_n\in\mathbb Z[X].
\]
This is the polynomial conclusion before the source's subsequent
integer-specialisation step in display~(24).
The coefficientwise interpretation also follows by comparing the source
rational coefficients: $F$ is not a rational function, since
$F(e^h)=h^{-1}\log(1/h)+O(h^{-1})$ as $h\downarrow0$.
For this estimate split the original sum at $n=1/h$; the remaining geometric
tail is $O(h^{-1})$.

In the source expression for $A_n$, the degree of the $k$th summand is
\[
 d_k=a_0k+E_k+(a_1-1)(k-a_1)+(\beta-k-1)(k-a_2),
\]
where
\[
 E_k=\frac{a_1(a_1-1)-(\beta-a_2)(\beta-a_2-1)
                   +(\beta-k)(\beta-k-1)}2.
\]
For $a_2\le k\le\beta-2$, one has $d_{k+1}-d_k=40n+1-k>0$.
There is therefore a unique highest-degree summand, and
\[
 K_n:=\deg A_n=\frac{1091n^2+81n+2}{2},\qquad
 W_n:=\deg U_n=K_n-M_n+\sum_{\ell\le15n}(1-\nu_\ell)\varphi(\ell).
\]
For fixed $n$, the positive representation below gives $H_n(x)=O(1)$ as
$x\to\infty$.  Hence $\Lambda_n(x)=O(x^{W_n-K_n})$.
Since $F(x)=O(x^{-1})$ and $K_n\ge1$, the identity
$V_n=U_nF-\Lambda_n$ gives $\deg V_n\le W_n-1$.
Thus $W_n$ clears both coordinates.

\paragraph{The limiting degree cost.}
The elementary summatory estimate
$\sum_{\ell\le y}\varphi(\ell)=3y^2/\pi^2+O(y\log y)$ gives
\[
 \frac1{n^2}\sum_{\ell\le15n}\varphi(\ell)\longrightarrow\frac{675}{\pi^2}.
\]
For $[u,v)\in\mathcal I$, the condition $\{n/\ell\}\in[u,v)$ is the
disjoint union of intervals
$n/(k+v)<\ell\le n/(k+u)$ for $k\ge0$.
For finitely many $k$ the same summatory estimate applies term by term.
The remaining terms have total normalised mass $O(1/K^2)$ after truncation
at $k=K$, since they involve only $\ell\le n/(K+u)$.
Letting first $n$ and then $K$ tend to infinity proves
\[
 \frac1{n^2}\sum_{\ell\le15n}\nu_\ell\varphi(\ell)
 \longrightarrow\frac3{\pi^2}
 \sum_{[u,v)\in\mathcal I}\sum_{k\ge0}
 \left(\frac1{(k+u)^2}-\frac1{(k+v)^2}\right)=\frac{3J}{\pi^2}.
\]
Here no contributing index exceeds $14n$, since $u\ge1/14$.
Consequently
\[
 K_n/n^2\longrightarrow C_1,\qquad
 (K_n-W_n)/n^2\longrightarrow C_0.
\]

\paragraph{A positive remainder at every fixed real base.}
For $x>1$ and $q=1/x$, the source identity has the positive representation
\[
 H_n(x)=\sum_{t\ge0}q^{a_0t}
 \frac{(q^{t+1};q)_{a_1-1}}{(q;q)_{a_1-1}}
 \frac{(q;q)_{\beta-a_2-1}}{(q^{a_2+t};q)_{\beta-a_2}}.
\]
The last denominator has length $\beta-a_2$, as required by the source's
gamma expression~(7) and residues~(8).  The shorter length in one unnumbered
product on printed p.~156 is an indexing discrepancy.
Every finite product lies between $P=(q;q)_\infty>0$ and $1$, so
\[
 P^2\le H_n(x)\le\frac{P^{-2}}{1-q^{a_0}}.
\]
In particular $\log H_n(x)=O_x(1)$ and $\Lambda_n(x)>0$.

For fixed $x>1$, the cyclotomic identity
\[
 \log\Phi_\ell(x)-\varphi(\ell)\log x
 =\sum_{d\mid\ell}\mu(d)\log(1-x^{-\ell/d})
\]
has total absolute error $O_x(n)$ over $\ell\le15n$.
Indeed, it is at most $15n B(x)$, where
\[
 B(x)=\sum_{d\ge1}\frac{-\log(1-x^{-d})}{d}<\infty.
\]
Since $0\le1-\nu_\ell\le1$, it follows that
\[
 \log\Lambda_n(x)=-(K_n-W_n)\log x+O_x(n).
\]

\paragraph{Homogenisation.}
Because $U_n,V_n$ are integral polynomials of degree at most $W_n$,
$b^{W_n}U_n(a/b)$ and $b^{W_n}V_n(a/b)$ are integers.
That integrality and the cleared linear-form identity it produces are
\href{https://github.com/wcook04/plectis-erdos/blob/f36a98bf3d3e6f65f1074e3b800e3293d5b8a51a/ErdosProblems/Erdos1049/PaperHomogenisationR7.lean#L61}{Lean source}.
Moreover
\[
\begin{split}
 \log\bigl(b^{W_n}\Lambda_n(a/b)\bigr)
 &=K_n\log b-(K_n-W_n)\log a+O_{a/b}(n)\\
 &=\bigl(C_1\log b-C_0\log a\bigr)n^2+o(n^2).
\end{split}
\]
The coefficient is negative under the theorem's hypothesis.  Thus positive
integral linear forms in $F(a/b)$ tend to zero.  If $F(a/b)=r/s$ were rational,
every such form would have absolute value at least $1/|s|$, a contradiction.
The separation bound for a nonzero integral form at a rational target is
\href{https://github.com/wcook04/plectis-erdos/blob/f36a98bf3d3e6f65f1074e3b800e3293d5b8a51a/ErdosProblems/Erdos1049/TwoSelectorRemainderEscape.lean#L128}{Lean source}.
\end{proof}

\begin{corollary}\label{res:thirtyone-four}
$F\bigl((31/4)^r\bigr)$ is irrational for every integer $r\ge1$.
\end{corollary}
\begin{proof}
The exact inequalities $31^2<4^5$ and $4^{200}<31^{81}$ give
$2/5<\log4/\log31<81/200$.
The first term of each trigamma difference yields
\[
 J\ge J_0:=\sum_{[u,v)\in\mathcal I}(u^{-2}-v^{-2})
   =\frac{2015640690251}{25971865920}.
\]
Using $\pi>157/50$ gives the rational lower bound
\[
 \theta^*>\frac{2359630009523263}{5820307922172744}
           >\frac{81}{200}.
\]
Finally, the logarithmic ratio and coprimality are preserved by a common
positive integer power.
\end{proof}

\begin{corollary}[an irrationality measure uniform over powers]
\label{cor:rational-base-measure}
For coprime $a>b\ge1$ with $\theta=\log b/\log a<\theta^*$ and every
integer $r\ge1$,
\[
 \mu_{\rm irr}\!\left(F((a/b)^r)\right)
 \le\frac{1-\theta}{\theta^*-\theta}.
\]
Here $\mu_{\rm irr}(\xi)$ is the supremum of the exponents $\nu$ for which
$|\xi-p/q|<q^{-\nu}$ has infinitely many reduced rational solutions.
In particular, $\mu_{\rm irr}(F((31/4)^r))<301$ for every $r\ge1$.
\end{corollary}
\begin{proof}
The raw source coefficient is a sum of $O(n)$ Laurent monomials times two
Gaussian polynomials, each of coefficient sum at most $2^{27n+2}$.
Its coefficient norm is $\exp(O(n))$, and every exponent is at most $K_n$.
The same normalising multiplier used above therefore gives
\[
 |U_n(x)|\le x^{W_n}\exp(O_x(n))\qquad(x>1\text{ fixed}).
\]
Set $\xi=F(a/b)$, $Q_n=b^{W_n}U_n(a/b)$ and
$P_n=b^{W_n}V_n(a/b)$.  With
\[
 \alpha=(C_1-C_0)\log a,\qquad
 \tau=C_0\log a-C_1\log b>0,
\]
we have $\log(Q_n\xi-P_n)=-\tau n^2+o(n^2)$ and
$|Q_n|\le\exp(\alpha n^2+o(n^2))$.

For integers $A,B,p,q$ with $q>0$, if $L=A\xi-B$ and $2q|L|\le1$, then
\[
 |L|\le |A|\,|\xi-p/q|.
\]
When $Ap-Bq=0$ this is equality.  Otherwise the nonzero integer $Ap-Bq$
gives $1/q\le |L|+|A|\,|\xi-p/q|$, which proves the inequality.
For a fixed $\eta\in(0,\tau)$ choose
$n=\lceil\sqrt{\log(2q)/(\tau-\eta)}\rceil$ and apply it to $(Q_n,P_n)$.
The two-sided remainder estimate yields
\[
 |\xi-p/q|\ge
 q^{-(\alpha+\tau+2\eta)/(\tau-\eta)-o(1)}.
\]
Let $\eta\downarrow0$.  The resulting bound is
$1+\alpha/\tau=(1-\theta)/(\theta^*-\theta)$.
Taking a common power multiplies $\alpha,\tau$ by $r$, leaving this quotient
unchanged; the constants in the approximation inequality may depend on $r$.
For $31/4$, sixteen terms of each trigamma difference and
$\pi>314159/100000$ give $\theta^*>0.40568$.
The logarithm series gives $\log4/\log31<0.4036982$.
Thus the bound is less than $2981509/9909<301$.
\end{proof}

\paragraph{Comparison and scope.}
Bundschuh and V\"a\"an\"anen's Theorem~2 at $\alpha=-1$
\cite[p.~177]{bv1994} gives
$\log b/\log a<\theta_{\rm BV}:=1/2-1/\pi^2$.
Since $\pi^2<10$, one has $\theta_{\rm BV}<2/5<\log4/\log31$.
The displayed sufficient regions therefore differ on
$[\theta_{\rm BV},\theta^*)$.
The direction $(14,12,14;27)$, its thirteen intervals and the constant
$\mu\approx2.46497868$ are inherited from \cite[p.~162]{zudilin2004},
where $\mu$ bounds an integer-base irrationality exponent.
The rational-base extension announced in \cite[Section~2]{zudilin2016}
has the same shape with an unspecified computable constant.
Here the polynomial specialisation identifies $\mu$ as admissible for $F$;
no priority claim is attached.

Negative bases are not treated.

\subsection{The base-uniform degree restriction}
The preceding construction uses degree and decay estimates valid at every
fixed real base greater than one.  Such uniformity itself imposes a limit.

\begin{theorem}[Archimedean cap on base-uniform rank-two families]
\label{res:archimedean-cap}
Let $(U_n,V_n)$ be pairs in $\mathbb Z[X]^2$ satisfying $\Lambda_n(x)=U_n(x)F(x)-V_n(x)\ne0$,
$\deg U_n,\deg V_n\le\delta n^2(1+o(1))$,
$\log\max(H(U_n),H(V_n))\le h n^2(1+o(1))$ with $H$ the $\ell^1$ coefficient
norm, and
$\log|\Lambda_n(x)|=-\sigma n^2\log x\,(1+o(1))$ for every real $x>1$, with
$\sigma,\delta>0$ and $h\ge0$ independent of $x$.  Then with
$d_n=\max(\deg U_n,\deg V_n)$, the homogenised forms
$b^{d_n}\Lambda_n(a/b)$ tend to zero whenever
$\log b/\log a<\sigma/(\sigma+\delta)$, and $\sigma/(\sigma+\delta)\le1/2$.
\end{theorem}

\begin{proof}
Write $H_n=\max(H(U_n),H(V_n))$.  Then
$|U_n(x)|,|V_n(x)|\le H_n x^{d_n}$ for real $x>1$.
Suppose $\sigma>\delta$ and choose an integer $p\ge2$ with
$(\sigma-\delta)\log p>h$.  Set $a_n=U_n(p)$, $b_n=V_n(p)$ and
$L_n=a_n F(p)-b_n$.  Hypotheses on height and degree give
$|a_n|\le\exp((h+\delta\log p)n^2+o(n^2))$, while
$|L_n|=\exp(-\sigma\log p\,n^2+o(n^2))$.  The adjacent integer
\[
 a_nb_{n+1}-a_{n+1}b_n=a_{n+1}L_n-a_n L_{n+1}
\]
is then $o(1)$, hence eventually zero.  Also $a_n\ne0$ for large $n$:
otherwise the nonzero integer $L_n=-b_n$ would have absolute value less
than $1$.  Thus $b_n/a_n$ is eventually a fixed rational $r$.  If
$F(p)\ne r$ then $|L_n|\ge|F(p)-r|$; if $F(p)=r$ then $L_n=0$.  Both
contradict the hypotheses, so $\sigma\le\delta$.
The integer argument of this paragraph, from the two cross-product limits to
the contradiction, is \href{https://github.com/wcook04/plectis-erdos/blob/f36a98bf3d3e6f65f1074e3b800e3293d5b8a51a/ErdosProblems/Erdos1049/PaperRankTwoCapR7.lean#L123}{Lean source}.
The homogenised logarithm satisfies
\[
 \limsup_{n\to\infty} n^{-2}\log\bigl|b^{d_n}\Lambda_n(a/b)\bigr|
 \le \delta\log b-\sigma\log(a/b),
\]
which is negative on the stated sufficient region.
The arithmetic form of that region is
\href{https://github.com/wcook04/plectis-erdos/blob/f36a98bf3d3e6f65f1074e3b800e3293d5b8a51a/ErdosProblems/Erdos1049/PaperRankTwoCapR7.lean#L144}{Lean source}, and the final numerical clause is
\href{https://github.com/wcook04/plectis-erdos/blob/f36a98bf3d3e6f65f1074e3b800e3293d5b8a51a/ErdosProblems/Erdos1049/PaperRankTwoCapR7.lean#L155}{Lean source}.
The conclusion bounds the
sufficient cutoff furnished by the displayed degree estimate.  An exclusion for a
particular family requires its actual degree and remainder asymptotics.
\end{proof}

\section{Exact normalized-Hankel order in Zudilin's construction}
\label{sec:hankel-order}

At $x=z=1$, Zudilin's normalized moments are
\[
 v_m^*=\sum_{t\ge0}q^{(m+1)t}
 \frac{(q;q)_m^3(q^{t+1};q)_m}{(q^{m+1+t};q)_{m+1}},
 \qquad
 V_N^*=\det_{0\le i,j<N}(v_{i+j}^*).
\]
His row transformation proves
\[
 \operatorname{ord}_q V_N^*\ge
 \frac{N(N-1)(2N-1)}6
\]
for every $N\ge1$~\cite[Section~4]{zudilin2016}.  A formal moment expansion identifies the unique least-order term and
shows that this estimate is always sharp.

\begin{theorem}[sharp normalized Hankel order]
\label{res:zudilin-sharp-qorder}
For every $N\ge1$,
\[
 \operatorname{ord}_q V_N^*=\frac{N(N-1)(2N-1)}6,
\]
and the coefficient of the first nonzero monomial is
\[
 [q^{N(N-1)(2N-1)/6}]V_N^*
   =\frac{(N!)^2(N+1)!}{2^N}.
\]
\end{theorem}

\begin{proof}
Write $P=(q;q)_\infty$ and introduce the coefficientwise formal series
\[
 G_q(w)=\frac1{(w;q)_\infty^3}
 \sum_{t\ge0}\frac{w^t}{(q;q)_t}
       \frac{(q^tw^2;q)_\infty}{(q^tw;q)_\infty^2},
 \qquad a_k(q)=[w^k]P^4G_q(w).
\]
For fixed $w$-degree, all operations define power series in $q$.
The product identities
$(q;q)_m=P/(q^{m+1};q)_\infty$ and
$(q^a;q)_m=(q^a;q)_\infty/(q^{a+m};q)_\infty$
give the exact formal identity
\[
 v_m^*=P^4G_q(q^{m+1})
      =\sum_{k\ge0}a_k(q)q^{(m+1)k}.
\]
At $q=0$, the term $t=0$ in $G_q$ is
$(1-w^2)/(1-w)^5$ and the terms $t\ge1$ sum to $w/(1-w)^4$.
Consequently
\[
 G_0(w)=\frac{1+2w}{(1-w)^4},\qquad
 a_k(0)=\frac{(k+1)^2(k+2)}2=:c_k>0.
\]

Apply Cauchy--Binet to a finite truncation of the moment sum and then pass
coefficientwise to the limit.  This is legitimate because only finitely many
increasing index tuples contribute to any fixed $q$-degree.  It gives
\[
 V_N^*=\sum_{k_0<\cdots<k_{N-1}}
 \left(\prod_{i=0}^{N-1}a_{k_i}(q)q^{k_i}\right)
 \prod_{0\le i<j<N}(q^{k_i}-q^{k_j})^2.
\]
The summand indexed by $(k_i)$ has order
\[
 \sum_{i=0}^{N-1}(2N-1-2i)k_i.
\]
Every weight is positive and $k_i\ge i$.  Equality with the least possible
order occurs uniquely when $k_i=i$ for every $i$.  That tuple contributes
\[
 \sum_{i=0}^{N-1}(2N-1-2i)i=\frac{N(N-1)(2N-1)}6
\]
and leading coefficient
\[
 \prod_{i=0}^{N-1}c_i=\frac{(N!)^2(N+1)!}{2^N}.
\]
No other tuple can cancel this coefficient.
\end{proof}

\paragraph{Formal order and fixed-base size.}
The theorem identifies the first formal term.  Formal order alone would not
control a fixed-$q$ residual: multiplying by $(1-q)^{N^3}$ preserves that first
term and changes the logarithm by a cubic quantity at fixed $0<q<1$.
The separate positive-measure argument for these same moments proves
$V_N^*(q)>0$ and
$\log(V_N^*(q)/(C_Nq^{B_N}))=O_q(N)$, where
$B_N=N(N-1)(2N-1)/6$ and $C_N=(N!)^2(N+1)!/2^N$.
Its complete ordinary proof is in
\texttt{ZudilinHankelPositiveMeasure.md} and the long record.
Neither fact supplies a denominator factor for the 2004 polynomial forms.



\section{Local cancellation and the remaining real estimate}
\label{sec:open}

\begin{theorem}[no corridor at base $3/2$]\label{res:nocorridor}
For all $N\ge1$ and $K\ge1$ and all natural $Q,D$, the tuple
$(3,2,N,K,Q,D)$ is not a coordinatewise corridor.
\end{theorem}

\begin{theorem}[cleared-tail recurrence]\label{res:tailrec}
Let $r,s,B,F\in\Q$ with $r\ne0$, let $c:\N\to\Q$, and let $P_N$ and $U_N$ be
the cleared tail state.  Then for every $N$,
\[
 U_{N+1}=r\,U_N-B\,c(N+1)\,s^{\,N+1}.
\]
\end{theorem}

\begin{theorem}[the forcing term]\label{res:forcing}
Let $s,B$ be natural numbers and $c:\N\to\N$, and put
$G_N=B\,c(N+1)\,s^{\,N+1}$.
\begin{enumerate}
\item If $s\ge2$, $B\ge1$ and $c(N+1)\ge1$, then $2^{\,N+1}\le G_N$.
\item If $s=1$, then $G_N=B\,c(N+1)$.
\end{enumerate}
\end{theorem}

Fix a common width $W$.  For $P\in\Z[X]$ of degree at most $W$, write
\[
 H_W(P)=2^WP(3/2)=\sum_{i=0}^W[P]_i3^i2^{W-i}.
\]
A specialised row is primitive when its two integer coordinates have gcd
one.  This normalisation is different from removing the common polynomial
coefficient content, and must precede the local count.
For depths $R,S\ge0$, define
\[
 J_{3,R}(P)=H_W(P)\pmod{3^R},\qquad
 J_{2,S}(P)=H_W(P)\pmod{2^S}.
\]
All four jets of $(U,V)$ vanish precisely when $D=3^R2^S$ divides both
specialised coordinates.  The dependence on the declared width remains fixed
when rows are added.

\begin{theorem}[binary four-jet collision]\label{res:jetkernel}
Fix a width \(W\) and depths \(R,S\), and let \((U_j,V_j)_{j<M}\) be any \(M\)
pairs of integral polynomials.  Call a subset of \(\{0,\dots,M-1\}\), equivalently
a vector of \(\{0,1\}^M\), a \emph{binary selector}.  If the
\(2^M\) binary selectors outnumber the finite four-jet target
\[
 (\Z/3^R\Z)^2\times(\Z/2^S\Z)^2,
\]
then two distinct subsets have the same four-jet sum.  Subtracting their
indicator vectors gives a nonzero coefficient vector in
\(\{-1,0,1\}^M\) cancelling all four jets.  The target has exact cardinality
\[
 (3^R)^2(2^S)^2.
\]
In particular, if \(R>0\) and \(4R+2S\le M\), such a collision exists.
\end{theorem}

\begin{proof}
Send each binary selector to the sum of the four-jet signatures it selects.
The claimed cardinal inequality and the pigeonhole principle give two
distinct selectors in the same fibre.  The cardinality formula is the product
of the four cyclic-modulus cardinalities.  For \(R>0\),
\[
 (3^R)^2(2^S)^2<(4^R)^2(2^S)^2=2^{4R+2S}\le2^M,
\]
which proves the stated sufficient threshold.
The cardinality formula, the collision, the signed \(\{-1,0,1\}\) vector and
the sufficient width are together \href{https://github.com/wcook04/plectis-erdos/blob/f36a98bf3d3e6f65f1074e3b800e3293d5b8a51a/ErdosProblems/Erdos1049/PaperFiniteAssembliesR7.lean#L248}{Lean source}.
\end{proof}

The ambient count does not use relations between the two residue
coordinates.  Vanishing minors can reduce that cost to one coordinate.

\begin{theorem}[B\'ezout--Pl\"ucker tail collapse]
\label{res:plucker-collapse}
Let $R_0$ be a commutative ring and let $w_n=(A_n,B_n)\in R_0^2$.  Suppose
that every row is unimodular ($u_n A_n+v_n B_n=1$ for some $u_n,v_n$) and every adjacent minor vanishes:
\[
 A_nB_{n+1}-B_nA_{n+1}=0\qquad(n\ge0).
\]
Then every pairwise minor $A_iB_j-B_iA_j$ vanishes.  In particular, take
$R_0=\Z/(2^S3^R)\Z$ with $R>0$.  If $S+2R\le k$, there are two distinct
binary selectors $s,t\in\{0,1\}^{k}$ such that
\[
 \sum_{i<k}s_iw_i=\sum_{i<k}t_iw_i.
\]
Thus the sufficient width is $S+2R$, rather than the ambient two-coordinate
width $2S+4R$.
\end{theorem}

\begin{proof}
A B\'ezout identity makes each row a unimodular anchor; a unit coordinate is the special case already recorded.  Vanishing of the
next minor therefore writes the next row as a scalar multiple of the current
one; induction places the entire tail on the line through $w_0$ and proves the
pairwise-minor assertion.  A determinant-one B\'ezout shear sends $w_0$ to
$(1,0)$, so all selector sums have only one free residue coordinate and occupy
at most $2^S3^R$ values.  Finally
\[
 2^S3^R<2^S4^R=2^{S+2R}\le2^k,
\]
and pigeonhole gives the two selectors.
The ring-generic minor collapse and the modular selector collision are
together \href{https://github.com/wcook04/plectis-erdos/blob/f36a98bf3d3e6f65f1074e3b800e3293d5b8a51a/ErdosProblems/Erdos1049/PaperFiniteAssembliesR7.lean#L269}{Lean source}.
\end{proof}

Primitive integer rows are unimodular modulo every modulus.  A particular
coordinate need not be a unit: $(2,3)$ modulo $6$ is an example.  The
minor-vanishing hypothesis remains a separate condition on the source rows.

\subsection{One minor gcd controls two different costs}

Let a rank-two lattice $\Lambda\subset\Z^2$ be generated by primitive rows,
and let $g>0$ be the gcd of their $2\times2$ minors.  Its Smith invariants are
$1,g$.  Therefore
\[
 \bigl|\operatorname{im}(\Lambda\bmod D)\bigr|
       =\frac{D^2}{\gcd(g,D)},\qquad
 \left[\Z^2:\frac{\Lambda\cap D\Z^2}{D}\right]
       =\frac{g}{\gcd(g,D)}.
\]
To verify the formulas, make a unimodular change of coordinates sending
$\Lambda$ to $\Z\oplus g\Z$; that change preserves $D\Z^2$.
The first quotient counts the modular signatures.  The second is the index
remaining after dividing a successful collision by $D$.
Once $D\mid g$, the modular count is $D$, while the remaining index is $g/D$.

If two independent divided rows $(A_i,B_i)$ have $|A_i|\le H$ and
$|A_i\xi-B_i|\le\varepsilon$, their determinant gives
\[
 2H\varepsilon\ge \frac{g}{\gcd(g,D)}.
\]
This is a restriction on two independent forms in the divided lattice.
For one nonzero integral form tending to zero, no additional product
$H\varepsilon\to0$ is required.

\subsection{A collision must have a small nonzero real remainder}

A real bin records the analytic requirement alongside the modular signature.
The multiplicity to control is the number of equal real values within one
modular fibre.

\begin{theorem}[quantitative bounded-fibre escape]\label{res:boundedfibre}
Let $A,B,J$ be finite sets and let $f:A\to B$, $g:A\to\mathbb R$ and
$\iota:A\to J$.  Suppose that each simultaneous fibre of $(f,g)$ has at
most $k$ elements and that, for some $\delta>0$,
\[
 \iota(x)=\iota(y)\quad\Longrightarrow\quad |g(x)-g(y)|<\delta.
\]
If $|B||J|k<|A|$, then some distinct $x,y\in A$ satisfy
\[
 f(x)=f(y),\qquad 0<|g(x)-g(y)|<\delta.
\]
\end{theorem}
\begin{proof}
Partition $A$ by $(f,\iota)$.  Some cell has more than $k$ elements, so its
$g$-values cannot all agree.  Two unequal values in that cell have the same
modular signature and differ by less than $\delta$.
\end{proof}

For primitive integer rows $(A_j,B_j)$, put $e_j=A_jF(3/2)-B_j$ and apply
the theorem to binary selectors, with
\[
 f(\varepsilon)=\sum_j\varepsilon_j(A_j,B_j)\pmod D,
 \qquad g(\varepsilon)=\sum_j\varepsilon_j e_j.
\]
Let $Q$ be the number of attained modular signatures and let $k$ bound the
simultaneous $(f,g)$ fibres.  All selector remainders lie in an interval of
length $T=\sum_j|e_j|$.  Dividing that interval into half-open bins of width
$D/n$ gives the sufficient inequality
\begin{equation}\label{eq:quantitative-selector-budget}
 2^M>Qk\left(\left\lfloor\frac{nT}{D}\right\rfloor+1\right).
\end{equation}
Its conclusion is a signed row sum divisible coordinatewise by $D$, with a
nonzero divided real remainder of absolute value less than $1/n$.
More precisely, if modular fibre $b$ has its own real span $T_b$ and exact
value multiplicity $k_b$, it suffices that
\[
 2^M>\sum_b k_b\left(\left\lfloor\frac{nT_b}{D}\right\rfloor+1\right).
\]
These are estimates for the primitive real remainders of the declared source
family.  Bounded unnormalised hypergeometric remainders cannot replace them.
A nonzero function need not be nonzero at $3/2$, as the factor $2X-3$ shows;
positive individual remainders need not remain positive after subtraction.

\subsection{The source-specific endpoint question}

The following sufficient construction keeps source membership, primitive
scaling, local cancellation and the real estimate together.  Without source
and height restrictions, arbitrary polynomial lifts of small rational
approximations would simply restate irrationality.

\begin{problem}[common-width simultaneous endpoint-jet construction]\label{prob:kernel}
Exhibit an integer constant \(C\ge1\) and, for every sufficiently large
positive integer \(n\), positive integers \(W_n,R_n,S_n,M_n\) such that
\[
 n^2\le W_n,R_n,S_n\le Cn^2,
 \qquad 4R_n+2S_n\le M_n\le Cn^2,
\]
together with polynomial pairs coming from a named source family with an
explicit coefficient-height bound after primitive normalisation,
\((U_{n,j},V_{n,j})\in\Z[X]^2\) for \(0\le j<M_n\), each of degree at most
the common declared width \(W_n\), whose specialised integer rows are
primitive:
\[
 \gcd\!\bigl(H_{W_n}(U_{n,j}),H_{W_n}(V_{n,j})\bigr)=1.
\]
Find a nonzero vector
\(\lambda^{(n)}\in\{-1,0,1\}^{M_n}\) for which, on putting
\[
 U_n=\sum_{j<M_n}\lambda^{(n)}_jU_{n,j},
 \qquad V_n=\sum_{j<M_n}\lambda^{(n)}_jV_{n,j},
\]
the pair \((U_n,V_n)\) is not \((0,0)\), all four common-width jets vanish,
\[
 J_{3,R_n}(U_n)=J_{3,R_n}(V_n)=0,
 \qquad J_{2,S_n}(U_n)=J_{2,S_n}(V_n)=0,
\]
where every jet in this display is formed using the declared width \(W_n\),
and the resulting divided integer linear form
\[
 A_n=\frac{H_{W_n}(U_n)}{3^{R_n}2^{S_n}},
 \qquad
 B_n=\frac{H_{W_n}(V_n)}{3^{R_n}2^{S_n}},
 \qquad
 \rho_n=A_nF(3/2)-B_n
\]
satisfies the explicit analytic condition
\[
 0<|\rho_n|<\frac1n.
\]
\end{problem}

The jet equations make $A_n,B_n$ integers.  If $F(3/2)=a/b$ were rational,
a nonzero $\rho_n$ would have absolute value at least $1/|b|$, contradicting
$|\rho_n|<1/n$ for large $n$.
The required analytic inequality is exactly the displayed scalar condition.
For irrationality it suffices to obtain such forms along any unbounded
sequence of $n$; the all-sufficiently-large-$n$ formulation is a stronger
construction requirement.  Height estimates become relevant when a specific
construction uses them to obtain this scalar inequality, or when an
independent-row determinant is invoked.

For the literal two-parameter source deformations, the long record keeps
the rowwise primitive normalisation, the two local minor valuations and the
remaining real-error minimum together.  The narrow regular-scale family has
an ordinary determinant obstruction to bounded divided errors; that
obstruction does not exclude sparse scales or the wider deformation.
Thus the next estimate must concern the real remainders of the chosen family.
Increasing the number of modular collisions alone does not establish it.

\paragraph{Functional equations.}
The classification of Bell and Smertnig implies that
$L(z)=\sum_{n\ge1}\tau(n)z^n$ is not $k$-Mahler for any $k\ge2$
\cite[Introduction]{bellsmertnig2026}.
The earlier simultaneous-$2$/$3$ obstruction is consequently subsumed by
this known single-base result.  A construction using additional functions or
functional relations must specify those functions and its closure conditions;
the single-base statement is not an obstruction to every approximation method.

\subsection*{Statements and declarations}

\paragraph{Artefact and data availability.} The
\href{\sourceurl}{pinned formal-source revision} contains the Lean sources, the
fixed toolchain, and the library manifest used in the verification.  The ordinary proofs used here are printed with their hypotheses.

\paragraph{Funding and competing interests.} This work received no external
funding.  The author declares no competing interests.

\paragraph{Acknowledgements.} The problem numbering and status follow the
Erd\H{o}s Problems catalogue maintained by Thomas Bloom~\cite{erdosproblems}.


\appendix
\section{Guide to the formal sources}\label{app:index}

Each linked phrase opens its Lean declaration at the pinned source
revision~\commitshort.  The declarations of this note live in seven modules:
\texttt{RationalBaseLambert}, \texttt{QAperyDiagonalNonEquivalence},
\texttt{RationalPadeArithmetic}, \texttt{ZudilinConeArithmetic},
\texttt{ZudilinHeightRegion}, \texttt{HermitePadeNoGo}, and
\texttt{BezoutPluckerJets}.  The first contains
the corridor, cleared-tail recurrence, and elementary $7/2$ certificate; the
second checks the finite $n=0$ diagonal residual; the remaining four separate
the Pad\'e exponent arithmetic, endpoint arithmetic, logarithmic comparisons,
rectangular exponent model, and B\'ezout--Pl\"ucker tail collapse.  The link
coordinates are validated against that pinned revision, so they remain correct
as later work moves lines in the working tree.

% The linked source manifest is retained in the authored TeX for automated
% validation, and kept outside the printed note.
\iffalse

\begin{tabular}{@{}ll@{}}
informal statement & linked Lean source\\
Van Assche diagonal initial values & \lrefx{Erdos1049/QAperyDiagonalNonEquivalence.lean}{33}{vanAsscheDiagonal_initial_values}\\
exact $n=0$ residual factorisation & \lrefx{Erdos1049/QAperyDiagonalNonEquivalence.lean}{67}{vanAsscheQAperyResidualAtZero_eq}\\
nonzero $n=0$ residual for $p>1$ & \lrefx{Erdos1049/QAperyDiagonalNonEquivalence.lean}{94}{vanAsscheQAperyResidualAtZero_ne_zero}\\
corridor predicate & \lrefx{Erdos1049/RationalBaseLambert.lean}{113}{CoordinatewiseCorridor}\\
corridor bound & \lrefx{Erdos1049/RationalBaseLambert.lean}{121}{coordinatewiseCorridor_implies_pow_lt_linear}\\
exponential comparison & \lrefx{Erdos1049/RationalBaseLambert.lean}{142}{three_mul_lt_two_pow_succ}\\
no corridor at $3/2$ & \lrefx{Erdos1049/RationalBaseLambert.lean}{155}{threeHalves_no_coordinatewiseCorridor}\\
rational-base prefix & \lrefx{Erdos1049/RationalBaseLambert.lean}{168}{rationalBasePrefixQ}\\
prefix step & \lrefx{Erdos1049/RationalBaseLambert.lean}{173}{rationalBasePrefixQ_succ}\\
cleared tail state & \lrefx{Erdos1049/RationalBaseLambert.lean}{181}{rationalBaseClearedTailQ}\\
cleared-tail recurrence & \lrefx{Erdos1049/RationalBaseLambert.lean}{187}{rationalBaseClearedTailQ_succ}\\
forcing magnitude & \lrefx{Erdos1049/RationalBaseLambert.lean}{198}{rationalBaseForcingNat}\\
exponential lower bound & \lrefx{Erdos1049/RationalBaseLambert.lean}{204}{twoPow_le_rationalBaseForcingNat}\\
integer-base collapse & \lrefx{Erdos1049/RationalBaseLambert.lean}{218}{rationalBaseForcingNat_one}\\
integer power comparison & \lrefx{Erdos1049/RationalBaseLambert.lean}{32}{sevenHalves_power_certificate}\\
exact rational boundary & \lrefx{Erdos1049/RationalBaseLambert.lean}{36}{sevenHalves_rational_margin}\\
logarithmic ratio bound & \href{https://github.com/wcook04/plectis-erdos/blob/99f4bf47422abbd8757cbb22b50ba079d764d3a7/ErdosProblems/Erdos1049/RationalBaseLambert.lean#L48}{\leanlabel{sevenHalves_log_ratio_lt}}\\
$\pi$ reciprocal-square bound & \href{https://github.com/wcook04/plectis-erdos/blob/99f4bf47422abbd8757cbb22b50ba079d764d3a7/ErdosProblems/Erdos1049/RationalBaseLambert.lean#L64}{\leanlabel{one_div_pi_sq_lt_one_nine}}\\
strict source margin & \href{https://github.com/wcook04/plectis-erdos/blob/99f4bf47422abbd8757cbb22b50ba079d764d3a7/ErdosProblems/Erdos1049/RationalBaseLambert.lean#L73}{\leanlabel{sevenHalves_bundschuhVaananen_margin}}\\
complete $7/2$ height condition & \lrefx{Erdos1049/RationalBaseLambert.lean}{83}{sevenHalves_archimedean_height_condition}\\
doubled common exponent & \lrefx{Erdos1049/RationalPadeArithmetic.lean}{19}{rationalPadeDenExpTwice}\\
doubled summand exponent & \lrefx{Erdos1049/RationalPadeArithmetic.lean}{24}{rationalPadePSummandDenExpTwice}\\
summand exponent bound & \lrefx{Erdos1049/RationalPadeArithmetic.lean}{30}{rationalPadePSummandDenExpTwice_le}\\
doubled maximal exponent & \lrefx{Erdos1049/RationalPadeArithmetic.lean}{46}{rationalPadeQMaxDenExpTwice}\\
exact gap identity & \lrefx{Erdos1049/RationalPadeArithmetic.lean}{52}{rationalPadeQMaxDenExpTwice_gap}\\
maximal exponent bound & \lrefx{Erdos1049/RationalPadeArithmetic.lean}{60}{rationalPadeQMaxDenExpTwice_le}\\
Pad\'e error scaling & \lrefx{Erdos1049/RationalPadeArithmetic.lean}{84}{rationalPadeError_mul}\\
row-content determinant factorisation & \lrefx{Erdos1049/RationalPadeArithmetic.lean}{94}{rationalPadeExteriorDet_mul_contents}\\
absolute determinant scaling & \lrefx{Erdos1049/RationalPadeArithmetic.lean}{104}{natAbs_rationalPadeExteriorDet_mul_contents}\\
content-product divisibility & \lrefx{Erdos1049/RationalPadeArithmetic.lean}{116}{contentProduct_dvd_rationalPadeExteriorDet}\\
exterior determinant identity & \lrefx{Erdos1049/RationalPadeArithmetic.lean}{124}{rationalPadeExteriorDet_cast_eq}\\
homogeneous endpoint evaluation & \lrefx{Erdos1049/ZudilinConeArithmetic.lean}{98}{homEvalThreeTwo}\\
four-jet target cardinality & \lrefx{Erdos1049/ZudilinConeArithmetic.lean}{131}{fourJetSignature_card}\\
four-jet pigeonhole kernel & \lrefx{Erdos1049/ZudilinConeArithmetic.lean}{142}{exists_distinct_binary_selectors_same_fourJet}\\
rank--depth collision threshold & \lrefx{Erdos1049/ZudilinConeArithmetic.lean}{161}{exists_distinct_binary_selectors_same_fourJet_of_rank}\\
bottom-jet divisibility & \lrefx{Erdos1049/ZudilinConeArithmetic.lean}{191}{bottomJet3_eq_zero_iff_dvd}\\
top-jet divisibility & \lrefx{Erdos1049/ZudilinConeArithmetic.lean}{197}{topJet2_eq_zero_iff_dvd}\\
bottom endpoint modulo three & \lrefx{Erdos1049/ZudilinConeArithmetic.lean}{203}{homEvalThreeTwo_mod_three}\\
generic homogeneous evaluation & \lrefx{Erdos1049/ZudilinConeArithmetic.lean}{221}{homEval}\\
left endpoint congruence & \lrefx{Erdos1049/ZudilinConeArithmetic.lean}{226}{homEval_mod_left}\\
right endpoint congruence & \lrefx{Erdos1049/ZudilinConeArithmetic.lean}{238}{homEval_mod_right}\\
left endpoint coprimality & \lrefx{Erdos1049/ZudilinConeArithmetic.lean}{254}{homEval_isCoprime_left_of_const_unit}\\
right endpoint coprimality & \lrefx{Erdos1049/ZudilinConeArithmetic.lean}{270}{homEval_isCoprime_right_of_top_unit}\\
two-endpoint coprimality & \lrefx{Erdos1049/ZudilinConeArithmetic.lean}{286}{homEval_isCoprime_mul_of_endpoint_units}\\
cyclotomic endpoint coprimality (Proposition~3.6) & \lrefx{Erdos1049/ZudilinConeArithmetic.lean}{302}{cyclotomicHomEval_isCoprime_mul}\\
top endpoint modulo two & \lrefx{Erdos1049/ZudilinConeArithmetic.lean}{320}{homEvalThreeTwo_mod_two}\\
unit constant endpoint obstruction & \lrefx{Erdos1049/ZudilinConeArithmetic.lean}{338}{three_not_dvd_homEvalThreeTwo_of_const_unit}\\
unit top endpoint obstruction & \lrefx{Erdos1049/ZudilinConeArithmetic.lean}{356}{two_not_dvd_homEvalThreeTwo_of_top_unit}\\
common-multiplier exclusion & \lrefx{Erdos1049/ZudilinConeArithmetic.lean}{398}{commonMultiplier_not_two_not_three_of_endpoint_units}\\
irrationality-exponent inequality & \lrefx{Erdos1049/ZudilinConeArithmetic.lean}{416}{twice_le_of_irrationalityExponent_bounds}\\
three-halves scalar margin & \lrefx{Erdos1049/ZudilinConeArithmetic.lean}{428}{three_two_scalar_margin_neg}\\
$31/4$ logarithmic-region membership & \lrefx{Erdos1049/ZudilinHeightRegion.lean}{45}{thirtyoneFour_mem_zudilinHeightRegion}\\
positive-power membership & \lrefx{Erdos1049/ZudilinHeightRegion.lean}{59}{thirtyoneFour_power_mem_zudilinHeightRegion}\\
$31/4$ outside the earlier region & \lrefx{Erdos1049/ZudilinHeightRegion.lean}{91}{thirtyoneFour_outside_bundschuhVaananenHeightRegion}\\
$81/200$ below the $3/2$ log ratio & \lrefx{Erdos1049/ZudilinHeightRegion.lean}{105}{eightyOneTwoHundredths_lt_threeHalves_log_ratio}\\
$3/2$ outside the $81/200$ region & \lrefx{Erdos1049/ZudilinHeightRegion.lean}{122}{threeHalves_outside_zudilinHeightRegion}\\
rectangular exponent-model bound & \lrefx{Erdos1049/HermitePadeNoGo.lean}{103}{rectangular_hp_threshold_le_classical}\\
equality case for the exponent model & \lrefx{Erdos1049/HermitePadeNoGo.lean}{126}{rectangular_hp_threshold_eq_classical_iff}\\
\end{tabular}
\fi

\begin{thebibliography}{99}
\small
\raggedright
\setlength{\itemsep}{0pt}
\bibitem{erdos1988} P.~Erd\H{o}s, \emph{On the irrationality of certain series:
  problems and results}, in A.~Baker (ed.), \emph{New Advances in Transcendence
  Theory}, Cambridge UP, 1988, pp.~102--109,
  doi:\href{https://doi.org/10.1017/CBO9780511897184.009}{10.1017/CBO9780511897184.009}.
\bibitem{bv1994} P.~Bundschuh and K.~V\"a\"an\"anen,
  \href{https://numdam.org/item/CM_1994__91_2_175_0.pdf}{\emph{Arithmetical
  investigations of a certain infinite product}},
  Compositio Math. \textbf{91} (1994), no.~2, 175--199.
\bibitem{zudilin2004} W.~Zudilin,
  \href{https://geodesic.mathdoc.fr/articles/10.4064/aa111-2-4/}{\emph{Heine's
  basic transform and a permutation group for \(q\)-harmonic series}},
  Acta Arith. \textbf{111} (2004), no.~2, 153--164,
  doi:10.4064/aa111-2-4.
\bibitem{zudilin2016} W.~Zudilin,
  \href{https://arxiv.org/abs/1601.02688}{\emph{On the irrationality of
  generalized \(q\)-logarithm}}, arXiv:1601.02688; Res. Number Theory
  \textbf{2} (2016),
  doi:\href{https://doi.org/10.1007/s40993-016-0042-x}{10.1007/s40993-016-0042-x}.
  The remark that the results extend to non-integer \(p=r/s\), \(|p|>1\), under
  an assumption \(\log|r|>c\log|s|\) for a computable \(c>0\), is at the end of
  Section~2; no value of \(c\) is computed there, and the remark is made for the
  generalized \(q\)-logarithm of that paper.
\bibitem{erdosproblems} T.~F. Bloom,
  \href{https://www.erdosproblems.com/1049}{\emph{Erd\H{o}s Problem \#1049}},
  \texttt{erdosproblems.com/1049}, accessed 28 July 2026
  (page displays ``last edited 28 September 2025'').  The current record
  labels the problem open, cites \textup{[Er88c, p.~102]} and \textup{[Er48]},
  and explicitly describes its status as the website owner's present
  assessment rather than a literature-completeness guarantee.
\bibitem{bellsmertnig2026} J.~Bell and D.~Smertnig,
  \href{https://arxiv.org/abs/2603.23456}{\emph{Mahler series with
  multiplicative coefficient sequences}}, arXiv:2603.23456v1,
  24 March 2026.  The introduction explicitly includes the divisor and
  totient functions among the examples which are not $k$-Mahler for any
  $k\ge2$.
\end{thebibliography}

\companionpapercontext

\end{document}
